\documentclass[11pt]{article}

% ---------------------------------------------------------------- packages
\usepackage{lmodern}
\usepackage[T1]{fontenc}
\usepackage[utf8]{inputenc}
\usepackage[margin=1in]{geometry}
\usepackage{amsmath,amssymb}
\usepackage{booktabs,longtable,array}
\usepackage{graphicx}
\usepackage{xcolor}
\usepackage{microtype}
\usepackage{enumitem}
\usepackage[most]{tcolorbox}
\usepackage{titlesec}
\usepackage{listings}
\usepackage{fancyhdr}
\usepackage[colorlinks=true,linkcolor=accent,urlcolor=accent,citecolor=accent]{hyperref}

% break long code identifiers at -, _, /, . without inserting hyphens
\urlstyle{tt}
\def\UrlBreaks{\do\/\do\-\do\_\do\.}
\newcommand{\cmd}[1]{\nolinkurl{#1}}

% ---------------------------------------------------------------- palette
\definecolor{accent}{HTML}{0B5394}
\definecolor{accentdark}{HTML}{073763}
\definecolor{rulegray}{HTML}{B7B7B7}
\definecolor{codebg}{HTML}{F5F6F8}
\definecolor{codeframe}{HTML}{D9DCE1}
\definecolor{statusbg}{HTML}{FBE9E7}
\definecolor{statusframe}{HTML}{C0392B}
\definecolor{notebg}{HTML}{EAF2FB}

% ---------------------------------------------------------------- section styling
\titleformat{\section}{\Large\bfseries\color{accentdark}}{\thesection}{0.6em}{}
\titleformat{\subsection}{\large\bfseries\color{accent}}{\thesubsection}{0.6em}{}
\titleformat{\subsubsection}{\normalsize\bfseries\color{accentdark}}{\thesubsubsection}{0.6em}{}
\titlespacing*{\section}{0pt}{2.2ex plus 1ex minus .2ex}{1.2ex plus .2ex}

% ---------------------------------------------------------------- headers
\pagestyle{fancy}
\fancyhf{}
\renewcommand{\headrulewidth}{0.4pt}
\renewcommand{\headrule}{\hbox to\headwidth{\color{rulegray}\leaders\hrule height \headrulewidth\hfill}}
\fancyhead[L]{\small\color{accentdark}\textbf{RaoRocketSim}}
\fancyhead[R]{\small\color{accentdark}Technical Reference}
\fancyfoot[C]{\thepage}

% ---------------------------------------------------------------- code listings
\lstdefinestyle{shell}{
  basicstyle=\ttfamily\footnotesize,
  backgroundcolor=\color{codebg},
  frame=single, rulecolor=\color{codeframe},
  breaklines=true, breakatwhitespace=true,
  showstringspaces=false, columns=fullflexible,
  xleftmargin=6pt, xrightmargin=6pt, framesep=6pt, aboveskip=8pt, belowskip=6pt,
}
\lstdefinestyle{py}{
  style=shell, language=Python,
  keywordstyle=\color{accent}\bfseries,
  commentstyle=\color{gray}\itshape,
  stringstyle=\color{accentdark},
}
\lstset{style=shell}

% ---------------------------------------------------------------- callout boxes
\newtcolorbox{statusbox}{colback=statusbg,colframe=statusframe,
  boxrule=0.8pt,arc=2pt,left=8pt,right=8pt,top=6pt,bottom=6pt}
\newtcolorbox{notebox}{colback=notebg,colframe=accent,
  boxrule=0.6pt,arc=2pt,left=8pt,right=8pt,top=6pt,bottom=6pt}

\setlist[itemize]{leftmargin=1.4em,topsep=2pt,itemsep=1.5pt}
\setlist[enumerate]{leftmargin=1.6em,topsep=2pt,itemsep=1.5pt}

\newcommand{\code}[1]{\texttt{#1}}

% ================================================================= document
\begin{document}

% ----------------------------------------------------------------- title page
\begin{titlepage}
\centering
\vspace*{2.2cm}
{\color{rulegray}\rule{\textwidth}{0.8pt}}\\[1.0cm]
{\Huge\bfseries\color{accentdark} RaoRocketSim\par}
\vspace{0.5cm}
{\Large A Differentiable Rao Thrust-Optimized-Parabola\\[2pt]
Rocket-Nozzle Design Toolbox\par}
\vspace{0.5cm}
{\large\itshape Technical Reference --- Methodology, Mathematics, and Physics\par}
\vspace{0.8cm}
{\color{rulegray}\rule{\textwidth}{0.8pt}}\\[1.4cm]

\begin{minipage}{0.85\textwidth}\centering\small
This reference compiles the methodology, governing equations, software
architecture, validation status, and roadmap of the RaoRocketSim project.
Every symbol, ratio, and governing equation has been cross-checked against
the source code and the primary literature in \code{propulsion\_texts/}.
\end{minipage}

\vfill
\begin{statusbox}
\centering\small
\textbf{Engineering status.} Every generated contour is stamped
\code{hardware\_qualified = False}. Passing the repository's design gates
means a result cleared internal preliminary screening only. It does
\textbf{not} replace independent CFD, conjugate heat-transfer analysis,
structural FEA, combustion-stability work, material allowables,
manufacturing review, inspection, proof testing, or hot-fire qualification.
\end{statusbox}
\vspace{0.8cm}

{\large Ibrahim Shahid\par}
\vspace{0.2cm}
{\color{accentdark} Reference generated June 2026}\par
\end{titlepage}

% ----------------------------------------------------------------- TOC
\tableofcontents
\thispagestyle{fancy}
\clearpage

% ================================================================= 1
\section{Overview}

RaoRocketSim is a Python research and preliminary-design toolbox for
axisymmetric rocket-nozzle aerodynamics. It pairs a practical Rao
thrust-optimized parabolic (TOP) bell workflow with quasi-one-dimensional
performance estimation, low-order engineering screens, an axisymmetric
method-of-characteristics (MOC) solver, a finite-dimensional Rao variational
boundary-value problem (BVP), a vendored NASA/JHU reference-code port, and a
differentiable JAX backend with exact-Jacobian solves and gradient
sensitivities.

The \textbf{default, trusted} path is the chart-based Rao/TOP
quadratic-B\'ezier contour: it is deterministic, endpoint-exact, and
benchmarked against published geometry. The MOC and variational solvers are
\textbf{active research implementations} --- they expose residuals, topology,
reference comparisons, and explicit reliability metadata, but they are not
promoted to design-validated or hardware-qualified status.

\subsection{Capabilities and maturity}

\renewcommand{\arraystretch}{1.25}
\begin{longtable}{p{0.20\textwidth}p{0.53\textwidth}p{0.17\textwidth}}
\toprule
\textbf{Area} & \textbf{Current implementation} & \textbf{Status} \\
\midrule
\endhead
Rao/TOP geometry & Upstream + downstream throat arcs and a quadratic B\'ezier
bell using interpolated Rao-chart angles & \textbf{Trusted baseline} \\
Ideal performance & Constant-$\gamma$, calorically perfect, quasi-1D
isentropic flow; $C_F$, thrust, $I_{sp}$, $c^*$, $\dot m$, exit state &
Preliminary \\
Thermochemistry & Built-in combustion-product constants, or optional
RocketCEA chamber $\gamma$, $M_w$, $T_c$, $c^*$ & Preliminary
(constant-$\gamma$ flow) \\
Direct MOC optimization & Axisymmetric characteristic march coupled to a
monotone spline wall under SLSQP / Nelder--Mead & Experimental \\
Rao variational / MOC BVP & NASA-topology seed, Rao stationarity,
characteristic compatibility, mass + length closure, $D$-state continuity,
validity diagnostics & Experimental research path \\
Differentiable solver & JAX/Optimistix Levenberg--Marquardt residual solve,
differentiable kernel march, optional solved $\theta_B$, exact $C_F$
sensitivities & Research-grade \\
Wall pressure \& separation & Quasi-1D wall pressure; Summerfield,
Kalt--Badal, Schmucker screens & Screening only \\
Thermal / cooling / structure & Boundary-layer displacement, Bartz-style heat
flux, regen-channel, thin-wall hoop-stress screens & Screening only \\
Geometry export & CSV, revolved STL, STEP (CadQuery B-rep or faceted AP214),
Inventor IPT manifest & CAD-review input \\
Validation & Unit/regression tests, literature manifests, NASA/JHU parsers,
parity checks & Software verification \\
\bottomrule
\end{longtable}
\renewcommand{\arraystretch}{1.0}

\begin{sloppypar}
The maturity column is enforced in code: each contour carries a
\cmd{design_status} keyed by method
(\cmd{preliminary_top_geometry}, \cmd{experimental_moc_geometry},
\cmd{experimental_variational_geometry},
\cmd{experimental_rao_variational_moc_bvp}) and a research-grade
\code{ContourReliability} level (\S\ref{sec:reliability}).
\end{sloppypar}

\subsection{Test and validation status}

The most recently recorded run of the normal (non-\code{slow}) selection, on
\textbf{June 14, 2026}, reports:

\begin{lstlisting}
784 passed, 4 xfailed, 26 deselected in 350.73 s
\end{lstlisting}

produced with

\begin{lstlisting}[language=bash]
MPLCONFIGDIR=/tmp/raosim-mpl \
  .venv-jax/bin/python -m pytest -q -m "not slow"
\end{lstlisting}

The repository contains 339 test functions, expanded by parametrization to
the counts above. The four expected \code{xfail}s record known research gaps:
the unresolved provenance of one historical NASA $TT'$ fixture, an unpackaged
Cuffel--Back--Massier dataset, and literature-promotion tests for the
experimental MOC and legacy variational paths. The 26 deselected tests are
long JAX solves, convergence studies, NASA fixed-end closure, and the full
Rao-chart sweep. This is \emph{software and mathematical verification}, not
physical validation against experiment or CFD.

% ================================================================= 2
\section{What the tool does}

Given a chamber pressure, ambient pressure, throat size or target thrust,
expansion ratio, propellant model, and contour method, RaoRocketSim can:

\begin{itemize}
\item generate a reduced-length Rao/TOP bell, a conical nozzle, an
MOC-optimized bell, or an experimental Rao variational/MOC contour;
\item size throat radius from a requested thrust, or size expansion ratio for
matched ideal expansion;
\item compute ideal exit Mach number and pressure, thrust coefficient,
thrust, mass flow, effective exhaust velocity, and specific impulse;
\item estimate wall pressure, overexpansion separation, altitude performance,
boundary-layer displacement, heat flux, regenerative-cooling capacity, and
thin-wall pressure stress;
\item generate simplified chamber and convergent geometry from $L^*$ and
contraction ratio;
\item sweep $\varepsilon$, $P_c$, or $R_t$; compare contour families; and run
literature-backed benchmark cases with explicit pass / report / xfail
policies;
\item plot contours, characteristic nets, Mach / pressure / angle fields,
wall distributions, exit-plane profiles, topology, residual diagnostics, and
JAX sensitivity fields;
\item export versioned CSV, STL, STEP, JSON, Markdown, and metadata
artifacts;
\item compute exact derivatives of the control-surface $C_F$ with respect to
the solved node variables via the JAX backend.
\end{itemize}

% ================================================================= 3
\section{Physical model}

\subsection{Assumptions}

The core gas-dynamics and nozzle solvers assume steady, inviscid, adiabatic
flow; a calorically perfect ideal gas with constant $\gamma$; isentropic
expansion except where an empirical loss or separation screen is applied;
axisymmetry for the active MOC and Rao implementations; and a choked throat
with a fully supersonic divergent section. There is \textbf{no} reacting-flow
chemistry evolution, finite-rate kinetics, particles, film cooling, wall
roughness, ablation, embedded shocks, side loads, or fluid--structure
interaction inside the solved flowfield. The optional CEA integration
supplies a chamber-property \emph{snapshot}; the nozzle flow is still
evaluated with a single effective chamber $\gamma$.

\subsection{Propellant / thermochemistry data}

The built-in database stores nominal \emph{combustion-product} properties
(the exhaust gas at the nominal mixture ratio), not raw propellant
properties:

\begin{center}
\begin{tabular}{lccccc}
\toprule
\textbf{Propellant} & $\gamma$ & $M_w$ [kg/mol] & $T_c$ [K] & $\eta_{Isp}$ & O/F \\
\midrule
N2O/Ethanol & 1.22 & 0.0260 & 2800 & 0.92 & 5.5 \\
LOX/RP-1    & 1.23 & 0.0235 & 3400 & 0.96 & 2.6 \\
LOX/LCH4    & 1.24 & 0.0220 & 3500 & 0.96 & 3.5 \\
LOX/LH2     & 1.20 & 0.0100 & 3250 & 0.98 & 6.0 \\
\bottomrule
\end{tabular}
\end{center}

The specific gas constant is $R = R_u / M_w$ with
$R_u = 8314.46~\mathrm{J\,kmol^{-1}K^{-1}}$. Users may also supply custom
properties or request RocketCEA-derived chamber values.

\subsection{Isentropic gas dynamics}

For Mach number $M$ and ratio of specific heats $\gamma$, the implemented
stagnation relations are
\begin{equation}
\frac{T}{T_0}=\left(1+\frac{\gamma-1}{2}M^2\right)^{-1},
\qquad
\frac{p}{p_0}=\left(\frac{T}{T_0}\right)^{\gamma/(\gamma-1)},
\qquad
\frac{\rho}{\rho_0}=\left(\frac{T}{T_0}\right)^{1/(\gamma-1)} .
\end{equation}
The area--Mach relation is
\begin{equation}
\frac{A}{A^*}=\frac{1}{M}
\left[\frac{2}{\gamma+1}\left(1+\frac{\gamma-1}{2}M^2\right)\right]^{\frac{\gamma+1}{2(\gamma-1)}},
\end{equation}
inverted with Newton iteration on the chosen (subsonic or supersonic) branch.
The Prandtl--Meyer function and Mach angle are
\begin{equation}
\nu(M)=\sqrt{\frac{\gamma+1}{\gamma-1}}\,
\tan^{-1}\!\sqrt{\frac{\gamma-1}{\gamma+1}\left(M^2-1\right)}
-\tan^{-1}\!\sqrt{M^2-1},
\qquad
\mu=\sin^{-1}\!\left(\frac{1}{M}\right),
\end{equation}
with $\nu(M)$ inverted by Newton iteration using
$\mathrm{d}\nu/\mathrm{d}M=\sqrt{M^2-1}\,/\,[M(1+\tfrac{\gamma-1}{2}M^2)]$.
These are standard compressible-flow results (Anderson, \emph{Modern
Compressible Flow}; cross-checked against \code{propulsion\_texts/prmeyer.pdf}).

\subsection{Performance model}

With $\varepsilon = A_e/A_t$, $A_t=\pi R_t^2$, and $A_e=\varepsilon A_t$, the
ideal one-dimensional thrust coefficient is
\begin{equation}
C_F=
\sqrt{\frac{2\gamma^2}{\gamma-1}
\left(\frac{2}{\gamma+1}\right)^{\frac{\gamma+1}{\gamma-1}}
\left[1-\left(\frac{p_e}{p_c}\right)^{\frac{\gamma-1}{\gamma}}\right]}
+\left(\frac{p_e-p_a}{p_c}\right)\varepsilon .
\end{equation}
The built-in propellant model applies an empirical efficiency multiplier
$\eta_{Isp}$ to the complete ideal coefficient:
\begin{equation}
C_{F,\mathrm{actual}}=\eta_{Isp}\,C_F,
\qquad
F=C_{F,\mathrm{actual}}\,p_c A_t .
\end{equation}
The characteristic velocity, mass flow, specific impulse, and effective
exhaust velocity are
\begin{equation}
c^*=\frac{\sqrt{\gamma R T_c}}{\gamma\sqrt{\left(2/(\gamma+1)\right)^{(\gamma+1)/(\gamma-1)}}},
\quad
\dot m=\frac{p_c A_t}{c^*},
\quad
I_{sp}=\frac{C_{F,\mathrm{actual}}\,c^*}{g_0},
\quad
V_e=I_{sp}\,g_0 ,
\end{equation}
with $g_0 = 9.80665~\mathrm{m\,s^{-2}}$.

\begin{notebox}
\small These are \textbf{ideal-cycle} estimates. The efficiency multiplier is
a single lumped factor, not a resolved loss model, and must not be read as a
prediction of combustion, boundary-layer, two-phase, or chemical losses.
(Formulation: Sutton \& Biblarz, \emph{Rocket Propulsion Elements};
Anderson.)
\end{notebox}

% ================================================================= 4
\section{Contour methods}

\subsection{Conical reference}

The conical utility reuses the bell throat arcs and a straight divergent wall
at half-angle $\alpha$. Its length and classical divergence factor are
\begin{equation}
L_{\mathrm{cone}}=\frac{R_e-R_t}{\tan\alpha},
\qquad
\eta_{\mathrm{div}}=\frac{1+\cos\alpha}{2} .
\end{equation}
The comparison module also estimates a bell divergence loss from an assumed
linear exit-plane flow-angle profile,
\begin{equation}
\eta_{\mathrm{div}}\approx
\frac{\displaystyle\int_0^{R_e}\rho u_x\,|\mathbf{u}|\,r\,\mathrm{d}r}
{\displaystyle\int_0^{R_e}\rho\,|\mathbf{u}|^2\,r\,\mathrm{d}r},
\qquad
C_{F,2D}\approx\eta_{\mathrm{div}}\,C_{F,1D},
\end{equation}
which is a comparison aid, not a resolved exit-plane solution of the B\'ezier
contour.

\subsection{Rao/TOP B\'ezier baseline (default, trusted)}

The baseline contour has three pieces: an upstream circular throat arc
($R_u = 1.5\,R_t$), a downstream circular throat arc ($R_d = 0.382\,R_t$),
and a quadratic B\'ezier bell from inflection point $N$ to exit point $E$.
The exit radius and reference $15^\circ$ cone length are
\begin{equation}
R_e=R_t\sqrt{\varepsilon},
\qquad
L_{15}=\frac{R_e-R_t}{\tan 15^\circ},
\qquad
L_n=\frac{L_{\%}}{100}\,L_{15},
\end{equation}
and the bell is
\begin{equation}
\mathbf{B}(t)=(1-t)^2\,\mathbf{N}+2(1-t)t\,\mathbf{P}_1+t^2\,\mathbf{E},
\qquad 0\le t\le 1,
\end{equation}
where $\mathbf{P}_1$ is the intersection of the tangent leaving $N$ at angle
$\theta_n$ and the tangent entering $E$ at angle $\theta_e$.

\paragraph{Chart angles and provenance.}
The default $(\theta_n,\theta_e)$ are bilinearly interpolated from embedded
Rao/NASA chart tables spanning approximately $4\le\varepsilon\le 50$ and
$60\%\le L_{\%}\le 100\%$. The tables reproduce Rao's TOP design charts
(Rao, \emph{Approximation of Optimum Thrust Nozzle Contour}, ARS J.\ 30(6),
1960; reproduced in Sutton and NASA SP-8120). Rao's underlying optimum study
was computed for $\gamma = 1.23$; per Rao (\emph{Recent Developments in
Rocket Nozzle Configurations}, ARS J.\ 31(11), 1961), the optimal
\emph{contour} is nearly $\gamma$-insensitive at fixed $(\varepsilon, L)$ ---
only $C_F$ depends strongly on $\gamma$ --- so the angle tables remain valid
comparison targets at other $\gamma$. This method is the trusted baseline
because it is deterministic, smooth, endpoint-exact, benchmarked against
explicit TOP geometry, and does not depend on the convergence of an
experimental flow solver.

\subsection{Chamber and convergent geometry}

Optional upstream geometry is sized from contraction ratio $CR=A_c/A_t$ and
characteristic length $L^*=V_c/A_t$:
\begin{equation}
R_c=R_t\sqrt{CR},
\qquad
V_c=L^* A_t .
\end{equation}
The code subtracts the convergent conical-frustum volume and assigns the
remainder to a cylindrical chamber. This is a geometric volume model only.

\subsection{Direct MOC wall optimization}

This path parameterizes the bell as a monotone cubic Hermite spline with
decision variables $\mathbf{q}=[\theta_n,\,r_1,\,r_2,\dots,r_{n_c}]$. For
each candidate wall it builds a transonic starting line, marches an
axisymmetric characteristic net with wall feedback, samples the exit plane,
and minimizes a cost combining negative exit thrust, exit-flow-angle
penalties, radius monotonicity, and curvature regularization. SciPy SLSQP is
used when available, with a NumPy Nelder--Mead fallback. The exported bell is
currently reconstructed as a smooth quadratic B\'ezier from the optimized
entrance and exit angles --- one reason it remains
\code{experimental\_moc\_geometry}.

\subsection{Rao variational / MOC BVP (main research solver)}

This solver seeds a finite-dimensional BVP from the NASA/JHU topology
\begin{equation}
T T' \;\rightarrow\; B \;\rightarrow\; BD \;\rightarrow\; D \;\rightarrow\; DE \;\rightarrow\; E .
\end{equation}
Its default configuration uses the JAX/Optimistix Levenberg--Marquardt
backend with exact autodiff Jacobians, the characteristic residual
formulation, a NASA-style fixed-end topology seed, full position /
flow-angle / Mach continuity at point $D$, and a continuation ladder that
raises weights on the mass, length, and endpoint constraints. The solved
control-surface unknowns are, schematically,
\begin{equation}
\mathbf{u}=\left[\,M_i,\;\theta_i,\;r_i \;\middle|\;
\lambda_2,\;\lambda_3,\;\log C,\;f_D\,\right],
\end{equation}
with optional wall unknowns and an optional live kernel angle $\theta_B$.
Here $f_D$ is the arc-length fraction locating $D$ on the kernel
characteristic $BD$.

\subsubsection{Axisymmetric characteristics}

With nodes ordered downstream and the axisymmetric source
$S=\sin\theta\,\sin\mu/r$, the implemented compatibility relations are
\begin{align}
C^{+}:&\quad \frac{\mathrm{d}r}{\mathrm{d}x}=\tan(\theta+\mu),
&& \mathrm{d}(\theta-\nu)=-S\,\mathrm{d}s, \\[2pt]
C^{-}:&\quad \frac{\mathrm{d}r}{\mathrm{d}x}=\tan(\theta-\mu),
&& \mathrm{d}(\theta+\nu)=+S\,\mathrm{d}s .
\end{align}
The $C^{+}$ control surface $DE$ is enforced geometrically by reconstructing
its axial coordinates,
\begin{equation}
x_{i+1}=x_i+\frac{r_{i+1}-r_i}{\tan(\bar\theta_i+\bar\mu_i)},
\qquad x_1=x_D(f_D) .
\end{equation}
The source term and $C^{\pm}$ invariant pairing are taken from Anderson
(\emph{Modern Compressible Flow} \S11.4) and Zucrow \& Hoffman (\emph{Gas
Dynamics} Vol.\ 2, Ch.\ 17), and are oracle-validated against the NASA
M3.5Perf grids (the corrected forms hold at RMS $\sim 10^{-6}$--$10^{-8}$ on
kernel and Deriv characteristics).

\subsubsection{Rao stationarity}

The critical Mach number is
\begin{equation}
M^*=\sqrt{\frac{(\gamma+1)M^2}{2+(\gamma-1)M^2}} .
\end{equation}
The primary algebraic optimum-thrust condition along the control surface $DE$
is
\begin{equation}
M^*\frac{\cos(\theta-\alpha)}{\cos\alpha}=C,
\qquad \alpha=\mu=\sin^{-1}(1/M),
\label{eq:rao1}
\end{equation}
implemented in logarithmic form for conditioning. This is Eq.~(1) of Rao,
Beck \& Booth, AIAA 99-2584 (\code{propulsion\_texts/rao1999.pdf}). Its
differential identity, Eq.~(2) of the same paper,
\begin{equation}
\mathrm{d}\ln M^*-(\mathrm{d}\theta-\mathrm{d}\alpha)\tan(\theta-\alpha)
+\mathrm{d}\alpha\tan\alpha=0,
\end{equation}
is retained as a secondary consistency check. The older direct-method
functionals are also available; per unit radial extent their
stagnation-normalized forms are
\begin{align}
f_1&=2\pi r\left[\left(\frac{p}{p_0}-\frac{p_a}{p_0}\right)
+\frac{\rho}{\rho_0}\gamma M^2\frac{T}{T_0}
\frac{\sin(\phi-\theta)\cos\theta}{\sin\phi}\right], \\
f_2&=2\pi r\frac{\rho}{\rho_0}\bar V\frac{\sin(\phi-\theta)}{\sin\phi},
\qquad
f_3=\cot\phi,
\end{align}
representing axial thrust, mass flow, and axial length respectively.

\subsubsection{Mass and endpoint closure}

Mass closure compares the surface-normal flux integral on $DE$ and the
selected $BD$ segment,
\begin{equation}
\dot m_\Gamma=\int_\Gamma 2\pi r\,\rho V\,\left|\sin(\beta-\theta)\right|\,\mathrm{d}s,
\qquad
\dot m_{DE}=\dot m_{BD} .
\end{equation}
The BVP additionally pins the commanded exit station and radius, $x_E=L_n$
and $r_E=R_t\sqrt{\varepsilon}$, and reports scaled residuals for mass,
length, stationarity, characteristic compatibility, geometry, regularization,
penalties, wall endpoints, and wall tangency.

\subsubsection{Smooth-flow validity}

The classical Rao validity inequality is evaluated along $DE$:
\begin{equation}
b=1-\frac{\mathrm{d}\alpha}{\mathrm{d}\theta}\,
\frac{\tan(\theta-\alpha)+\tan\alpha}{\tan(\theta-\alpha)-\tan\alpha}\;\ge\;0 .
\end{equation}
If $b$ drops below tolerance, the requested smooth variational solution lies
outside the shock-free Rao region; the construction is flagged
\code{invalid\_short\_nozzle\_region} and the reliability level is
downgraded. This boundary relation is the perfect-gas optimum-thrust boundary
of Rao, Beck \& Booth (1999) and is discussed in \"Ostlund's KTH thesis
(\code{propulsion\_texts/fulltext01.pdf}, \S3).

For the repository's $\varepsilon=10$, $80\%$-length, $\gamma=1.4$ regression
case, the smooth stationary-$DE$ reference is
\begin{lstlisting}
theta_B = 25.5659 deg
f_D     = 0.15216
D       = (M = 3.40145, theta = 18.5182 deg)
E       = (M = 3.47655, theta = 11.1193 deg)
\end{lstlisting}
a numerical regression point for the implemented formulation, not a universal
nozzle-design result.

\subsection{Transonic starting lines}

The MOC code provides several throat starting-line models: \code{kliegel\_levine}
(third-order toroidal-coordinate series; default), \code{sauer\_modified}
(compact leading-order curved-throat approximation), \code{area\_ratio}
(quasi-1D area--Mach line), \code{hall} (deprecated alias for
\code{sauer\_modified}), and \code{nasa\_visible\_kliegel\_levine} (a
source-faithful NASA/JHU compatibility mode used by the reference port). The
theory-correct Kliegel--Levine path documents and tests several
coefficient-transcription corrections; the source-faithful mode deliberately
preserves the historical C++ integer-division behavior required for NASA
binary/output parity.

% ================================================================= 5
\section{Differentiable (JAX) backend}

\code{raosim/jax/} provides the default inner solver for the Rao BVP and the
gradient API. The design keeps a single boundary: the NumPy shell in
\code{rao\_variational.solve\_rao\_bvp} still owns seeding, kernel
construction, reliability gating, diagnostics, and output assembly, while the
JAX layer owns only the inner least-squares solve.

\begin{itemize}
\item \textbf{Solver.} Optimistix Levenberg--Marquardt with the \emph{exact}
\code{jacfwd} Jacobian of the assembled residual, replacing SciPy's
finite-difference \code{least\_squares}. Box bounds are handled by a smooth
reparametrization $u = \mathrm{lo} + (\mathrm{hi}-\mathrm{lo})\,\sigma(z)$ so
every iterate stays strictly inside the box. A constraint-weight continuation
ladder up-weights the single-element mass/length/endpoint residuals; the
reported residual is always the unweighted one, so reliability gates stay
comparable across backends.
\item \textbf{Differentiable kernel march and $\theta_B$.} The NASA-style
kernel march and the live $\theta_B$ secant path are differentiable.
\item \textbf{Sensitivities.} \code{rao\_sensitivities(config, solution=...)}
returns exact derivatives of the control-surface $C_F$ with respect to the
solved node variables, explicit fixed-solution partials with respect to
$p_a/p_0$ and $\gamma$, and the residual-Jacobian condition number.
\end{itemize}

\textbf{Not yet provided (v2 deferrals):} total re-solved (implicit-function)
design derivatives with respect to $R_t$, $\varepsilon$, length percentage,
or $\gamma$; a differentiable sensitivity map carried through to the final
manufacturable bell wall; and the Hessian. The pinned, tested stack is
\code{jax==0.6.2}, \code{jaxlib==0.6.2}, \code{optimistix==0.0.11},
\code{equinox==0.13.8}.

% ================================================================= 6
\section{Screening models}

These models are deliberately low-order: useful for ranking concepts and
rejecting unsuitable designs, \textbf{not} for qualification.

\subsection{Wall pressure}
At each contour radius the code sets $A(x)/A_t=\left(r(x)/R_t\right)^2$,
inverts the area--Mach relation, and evaluates $p(x)/p_c$ isentropically. A
positive downstream pressure increment is flagged as non-monotonic (motivated
by NASA SP-8120).

\subsection{Separation}
Three empirical overexpansion-separation criteria are exposed (as
\emph{implemented}):
\begin{gather}
p_{\mathrm{sep}}\approx 0.4\,p_a \qquad\text{(Summerfield)}, \notag \\[2pt]
\frac{p_{\mathrm{sep}}}{p_a}\approx\frac{1}{1.88\,M_e-1}\qquad\text{(Kalt--Badal)}, \\[2pt]
\frac{p_{\mathrm{sep}}}{p_c}\approx\left(\frac{p_a}{p_c}\right)^{0.8}M_e^{-1}\qquad\text{(Schmucker)}. \notag
\end{gather}
The first contour station whose quasi-1D wall pressure falls below the chosen
threshold is reported as the estimated separation location. Side loads,
shock separation, hysteresis, and transient startup are not modeled.
(Criteria surveyed in NASA SP-8120 and Stark, 2009.)

\subsection{Boundary layer}
The displacement-thickness screen uses a turbulent flat-plate-style
correlation
\begin{equation}
\delta^*\approx 0.046\,\frac{s}{Re_s^{1/5}}\sqrt{\frac{T}{T_w}},
\end{equation}
and estimates an effective exit area ratio from $r_e-\delta_e^*$. Gas
viscosity uses a Sutherland-type estimate.

\subsection{Heat flux and cooling}
The heat-flux model preserves Bartz-like sensitivity to chamber pressure,
characteristic velocity, throat diameter, gas temperature, and wall
temperature, then applies empirical axial and area scaling
(\code{bartz\_style\_screening} --- not a full Bartz implementation). For
regenerative cooling, total absorbed heat and coolant rise are estimated by
\begin{equation}
\dot Q=\int q''(x)\,2\pi r(x)\,\mathrm{d}x,
\qquad
\Delta T_c=\frac{\dot Q}{\dot m_c\,c_{p,c}},
\end{equation}
followed by a one-resistance wall/convection temperature estimate. Channel
pressure drop, boiling, coking, critical heat flux, rib conduction, and
property variation are not solved.

\subsection{Structural screen}
The pressure-stress estimate is the thin-wall hoop relation
\begin{equation}
\sigma_h\approx\frac{p_c\,r_{\max}}{t_w}.
\end{equation}
Reported margins compare this stress, the estimated wall temperature, and the
peak heat flux against user-supplied material limits. Temperature-dependent
allowables, fatigue, creep, weld efficiency, buckling, and combined loads are
outside the model.

\subsection{Atmosphere and trajectory}
Altitude performance uses piecewise ISA layers (troposphere $-6.5$~K/km from
0--11~km, isothermal 11--20~km, $+1.0$~K/km 20--47~km) with an exponential
tail above 47~km. The optional trajectory module integrates a 1-D vertical
point mass,
\begin{equation}
m\,\dot v=F-\tfrac{1}{2}\rho v|v|\,C_D A-m\,g(h),
\qquad
g(h)=g_0\left(\frac{R_E}{R_E+h}\right)^2,
\end{equation}
with $R_E=6.371\times10^{6}~\mathrm{m}$. It is not wired into the main CLI and
does not model guidance, pitch, staging, throttling, winds, or 6-DOF
dynamics.

% ================================================================= 7
\section{Installation}

Python 3.12 is recommended because the differentiable backend is pinned to a
tested JAX stack.
\begin{lstlisting}[language=bash]
python3.12 -m venv .venv-jax
source .venv-jax/bin/activate
python -m pip install --upgrade pip
python -m pip install -r requirements.txt pytest
\end{lstlisting}
Core dependencies are NumPy, SciPy, Matplotlib, JAX, JAXlib, Optimistix, and
Equinox. Optional integrations (\code{rocketcea} for variable chamber
properties, \code{cadquery} for true revolved B-rep STEP) are not installed by
\code{requirements.txt}. Run commands from the repository root; the project
does not yet ship as an installable package.

% ================================================================= 8
\section{Command-line usage}

CLI pressure units are \textbf{bar} for $P_c$, \textbf{kPa} for $P_a$, and
\textbf{millimeters} for $R_t$ and manufacturing dimensions; internal APIs use
SI units.

\paragraph{Preliminary Rao/TOP design.}
\begin{lstlisting}[language=bash]
.venv-jax/bin/python main.py \
  --propellant LOX/RP-1 \
  --Pc 45 --Pa 101.325 \
  --Rt 20 --epsilon 10 --length-pct 80 \
  --method bezier --no-plot \
  --output nozzle_profile.csv
\end{lstlisting}
Output is written to a versioned directory \code{builds/vNNN\_YYYYMMDD\_HHMMSS/}.

\paragraph{Size the throat from thrust.}
\begin{lstlisting}[language=bash]
.venv-jax/bin/python main.py \
  --propellant LOX/LCH4 \
  --Pc 60 --target-thrust 10000 --epsilon 12 --no-plot
\end{lstlisting}
The sizing relation is
$R_t=\sqrt{F_{\mathrm{target}}/(\pi\,C_{F,\mathrm{actual}}\,p_c)}$.

\paragraph{Experimental Rao variational / MOC solve.}
\begin{lstlisting}[language=bash]
.venv-jax/bin/python main.py \
  --propellant LOX/RP-1 \
  --Pc 45 --Rt 20 --epsilon 10 \
  --method rao_variational_moc \
  --rao-moc-n-control 12 --rao-moc-n-kernel 12 \
  --rao-moc-max-nfev 200 --no-plot
\end{lstlisting}
This path can be expensive and remains experimental
(\cmd{--rao-moc-max-nfev} defaults to 25). \cmd{--rao-moc-skip-moc} runs a
faster residual-only study but prevents MOC reliability promotion.

\paragraph{Sweep and benchmark.}
\begin{lstlisting}[language=bash]
.venv-jax/bin/python main.py --propellant LOX/LCH4 \
  --Pc 60 --Rt 25 --epsilon 10 --sweep epsilon 4 50 20 --no-plot

.venv-jax/bin/python main.py \
  --benchmark-case lea_top_schomberg_2014 \
  --benchmark-method bezier \
  --benchmark-report builds/benchmarks --no-plot
\end{lstlisting}
Sweep variables are \code{epsilon}, \code{Pc}, \code{Rt}. Benchmark manifests
are \cmd{lea_top_schomberg_2014}, \cmd{rao_scarfed_moc_1990}, and
\cmd{vulcain_s1_separation_similarity}; each metric is classified
\code{strict}, \code{xfail}, or \code{report}.

% ================================================================= 9
\section{Python API}

\paragraph{Baseline contour and performance.}
\begin{lstlisting}[style=py]
from raosim.engine import compute_engine_performance
from raosim.nozzle_geometry import bell_nozzle_contour
from raosim.propellants import get_propellant

prop = get_propellant("LOX/RP-1")
contour = bell_nozzle_contour(
    Rt=0.020, epsilon=10.0, length_pct=80.0,
    method="bezier", gamma=prop.gamma,
)
performance = compute_engine_performance(
    Pc=4.5e6, Pa=101325.0, Rt=0.020, epsilon=10.0, prop=prop,
)
\end{lstlisting}

\paragraph{Design-gated workflow.}
\begin{lstlisting}[style=py]
from raosim.design import DesignInput, ThermoSpec, design_nozzle_v2

result = design_nozzle_v2(DesignInput(
    thermo=ThermoSpec(mode="constant_gamma", propellant_name="LOX/RP-1"),
    Pc=8.0e6, Rt=0.020, epsilon=8.0,
    method="bezier", mode="preliminary",
))
print(result.performance.thrust)
print(result.gate_report.to_dict())
\end{lstlisting}
The mode \code{validated} uses the stricter schema and gate set; it does
\textbf{not} set \code{hardware\_qualified = True}.

\paragraph{Rao BVP and sensitivities.}
\begin{lstlisting}[style=py]
from raosim.rao_variational import RaoSolverConfig, solve_rao_bvp
from raosim.jax import rao_sensitivities

config = RaoSolverConfig(
    Rt=0.020, epsilon=10.0, gamma=1.4, length_pct=80.0,
    solver_backend="jax", formulation="characteristic",
)
solution = solve_rao_bvp(config)
sens = rao_sensitivities(config, solution=solution)
print(solution.reliability, solution.residuals.max_scaled)
print(sens.cf, sens.condition_number)
\end{lstlisting}
Pass \code{solver\_backend="numpy"} to use the legacy SciPy path.

% ================================================================= 10
\section{Outputs}

Normal CLI runs create \code{builds/vNNN\_YYYYMMDD\_HHMMSS/} and may contain:
a contour CSV in meters; a binary STL inner surface or closed solid; a STEP
solid (CadQuery revolved B-rep when available, else a faceted AP214 fallback);
an Inventor conversion manifest pointing to the authoritative STEP file;
design-gate JSON and v2 physics-screening sections; sweep CSV or benchmark
JSON/Markdown reports; and a human-readable \code{metadata.txt}. Native
Autodesk Inventor IPT writing is not implemented; bolt patterns, injector
interfaces, throat inserts, tolerances, and weld/braze allowances are
metadata placeholders, not modeled solid features.

% ================================================================= 11
\section{Validation and reference evidence}

The repository deliberately separates \emph{software verification} from
\emph{physical validation}.

\subsection{Software and mathematical verification}
The test suite covers gas-dynamics identities and inverse relations; B\'ezier,
conical, chamber, curvature, export, and design-workflow behavior;
axisymmetric characteristic pairing and source terms; transonic
Kliegel--Levine coefficients and NASA-visible-source semantics; characteristic
topology, mass conservation, BDE closure, wall tangency, and crossing checks;
NumPy/JAX primitive and residual parity; exact-Jacobian convergence,
kernel-march parity, solved-$\theta_B$ derivatives, and $C_F$ sensitivities;
literature-manifest parsing with strict/report/xfail policies; NASA/JHU
legacy-output and Tecplot parsing; and plotting/diagnostic metadata.

\subsection{NASA/JHU reference port}
\code{Three-Dimensional-Nozzle-Design-Code-master/} vendors the Rice/JHU 2D
MOC, streamline-tracing, and 3D MOC source plus sample outputs. The active
Python port implements the initial throat line and kernel march; arc-wall,
interior, special-wall, and axis unit processes; mass flow along
right-running characteristics; point-$D$ selection, point-$E$ integration,
and fixed-/free-end closure; the $\theta_B$ secant solution; BDE-region and
wall-contour construction; and explicit \code{RaoTopology} objects for $TT'$,
$B$, $BF$, $D$, $BD$, $DE$, $E$, and the wall streamline.

The checked-in M3.5 perfect-nozzle wall comparison is within the repository's
$10^{-3}$ RMS regression gate (currently $\approx 1.8\times10^{-4}$ against
\code{wall.out}). However, the generator provenance of the historical
$TT'.\text{out}$ fixture is \textbf{unresolved}: the visible
\code{MOC\_GridCalc\_BDE.cpp} source-port workflow is the canonical reference
track, and the orphaned M3.5Perf fixture is a diagnostic overlay only --- not
promotion authority. Consequently the public solver does not label any
contour \code{NASA\_REFERENCE\_MATCHED}.

\subsection{Reliability labels}\label{sec:reliability}
The research code defines an ordered \code{ContourReliability} vocabulary:
\begin{center}
\code{geometric\_approximation} $\rightarrow$ \code{moc\_compatible}
$\rightarrow$ \code{rao\_variational\_residual\_solved} $\rightarrow$
\code{nasa\_reference\_matched} $\rightarrow$ \code{benchmark\_validated}
$\rightarrow$ \code{cfd\_checked} $\rightarrow$
\code{experimentally\_validated}.
\end{center}
This is a \emph{vocabulary}, not a claim that every level has been reached.
The \code{rao\_variational\_moc} solver assigns its level from per-solve
gates: a clean solve passing the BVP residual, MOC closure, valid-region, and
thrust-sanity gates reaches \code{rao\_variational\_residual\_solved}.
Promotion to \code{benchmark\_validated} is gated behind the module flag
\code{BENCHMARK\_VALIDATED\_AT\_RELEASE}, which is \textbf{\code{False}} in the
current code, so no public contour is benchmark-validated,
NASA-reference-matched, CFD-checked, or experimentally validated by this
repository.

% ================================================================= 12
\section{Repository layout}

\begin{lstlisting}
main.py                  Primary interactive/batch CLI
raosim/                  Active Python package
  gas_dynamics.py        Isentropic, area-Mach, Prandtl-Meyer relations
  engine.py              Ideal performance (Cf, Isp, c*, m_dot)
  propellants.py         Combustion-product property database
  nozzle_geometry.py     Rao/TOP Bezier geometry + chart tables
  conical.py             Conical reference contour + divergence factor
  chamber_geometry.py    L*/CR chamber + convergent sizing
  design.py              High-level schemas, gates, artifacts
  validation.py          Design-status labels and engineering gates
  physics.py             BL, thermal, cooling, structural screens
  separation.py          Summerfield / Kalt-Badal / Schmucker screens
  wall_pressure.py       Quasi-1D wall-pressure distribution
  atmosphere.py / altitude_performance.py / trajectory.py
  moc.py                 General axisymmetric MOC march + starting lines
  transonic_kernel.py    Kliegel-Levine transonic series
  nasa_moc.py            NASA/JHU-style kernel and BDE topology port
  moc_topology.py        RaoTopology objects (TT', B, BD, D, DE, E, wall)
  rao_optimizer.py       Direct MOC wall optimization
  rao_variational.py     Rao functionals, validity, and global BVP
  rao_residuals.py       Characteristic residual primitives
  jax/                   Differentiable primitives, residuals, kernel
                         march, LM solve, and sensitivities
  cea.py                 Optional RocketCEA thermochemistry
  benchmarks.py          Literature benchmark runner
  benchmark_data/        Manifests and digitized reference curves
  plotting.py            Geometry, field, topology, residual plots
  export.py              CSV, STL, STEP, IPT-manifest export
tests/                   Unit, regression, parity, benchmark tests
scripts/                 Research diagnostics and artifact generators
docs/                    NASA provenance, topology, and audit notes
latex-report/            Mathematical reference report and figures
propulsion_texts/        Local source literature used by the project
Three-Dimensional-Nozzle-Design-Code-master/
                         Vendored Rice/JHU reference source and outputs
\end{lstlisting}

% ================================================================= 13
\section{Known limitations}

\begin{enumerate}
\item \textbf{No hardware qualification.} All results require independent
analysis and test evidence.
\item \textbf{Constant-property nozzle flow.} CEA does not yet drive variable
$\gamma$, composition, or transport through the MOC field.
\item \textbf{Inviscid MOC.} Boundary-layer growth is a separate screen, not
coupled into marching or optimization.
\item \textbf{Incomplete separation physics.} Empirical onset criteria only.
\item \textbf{Experimental exact-Rao path.} Mathematically auditable and
tested, but release-level literature promotion is disabled
(\code{BENCHMARK\_VALIDATED\_AT\_RELEASE = False}).
\item \textbf{Chart vs.\ exact-variational angles.} Rao 1960 chart angles are
parabola-fit data; exact deltas are recorded, not forced to zero.
\item \textbf{Partial differentiability.} The complete
start-line-to-final-wall design map and total design derivatives are
unfinished.
\item \textbf{Low-order thermal/structural models.} Screens, not CHT,
coolant-network, fatigue, creep, or FEA.
\item \textbf{Simplified CAD.} Channels, bolt holes, injector faces, inserts,
and welds are not solid-modeled.
\item \textbf{No package metadata or stable API guarantee.}
\end{enumerate}

% ================================================================= 14
\section{Remaining work}

\begin{enumerate}
\item Complete and publish the exact-variational chart/reference sweep, define
accepted exact-vs-TOP deltas, and set criteria for flipping
\code{BENCHMARK\_VALIDATED\_AT\_RELEASE = True}.
\item Make the live differentiable $\theta_B$ solve routine and finish total
implicit design derivatives w.r.t.\ $R_t$, $\varepsilon$, $L_{\%}$, $\gamma$,
and ambient pressure.
\item Differentiate or replace the remaining start-line and BDE/final-wall
steps so sensitivities reach manufacturable wall coordinates.
\item Finish source-port provenance and public reliability wiring for the
NASA/JHU reference workflow, including the historical $TT'$ fixture.
\item Add variable-property frozen/equilibrium thermochemistry and transport
to the nozzle flowfield.
\item Couple viscous displacement, wall pressure, separation, and side-load
physics to the contour; validate against experiments and CFD.
\item Add independent axisymmetric Euler/RANS CFD comparisons with
quantitative gates for $C_F$, wall pressure, exit Mach, and flow angle.
\item Replace thermal/cooling/structural screens with conjugate heat transfer,
coolant models, temperature-dependent materials, and structural analysis.
\item Promote CAD from a revolved screening solid to a manufacturing
definition with channels, interfaces, fasteners, and tolerances.
\item Retire or isolate legacy paths, add package metadata, and publish a
stable API and reproducible release artifacts.
\end{enumerate}

% ================================================================= 15
\section{References}

Primary literature in \cmd{propulsion_texts/} and reference notes under
\code{docs/} and \cmd{raosim/benchmark_data/}:

\begin{itemize}
\item G. V. R. Rao, \emph{Exhaust Nozzle Contour for Optimum Thrust}, Jet
Propulsion, 1958.
\item G. V. R. Rao, \emph{Approximation of Optimum Thrust Nozzle Contour}, ARS
J.\ 30(6), 1960 --- the TOP $\theta_n/\theta_e$ design charts.
\item G. V. R. Rao, \emph{Recent Developments in Rocket Nozzle
Configurations}, ARS J.\ 31(11), 1961 --- $\gamma$-insensitivity of the
optimal contour.
\item Rao, Beck \& Booth, \emph{Rao Variational Optimum Bell Nozzle: A Design
Compendium}, AIAA 99-2584, 1999 --- stationarity Eqs.~(1)--(2), validity
boundary.
\item NASA SP-8120, \emph{Liquid Rocket Engine Nozzles} --- chart
reproduction, separation and wall-pressure guidance.
\item Rice, \emph{2D and 3D Method-of-Characteristics Tools for Complex Nozzle
Development}, JHU/APL RTDC-TPS-481, 2003.
\item Kliegel \& Levine, \emph{Transonic Flow in Small Throat Radius of
Curvature Nozzles}, 1969.
\item J. \"Ostlund, \emph{Flow Processes in Rocket Engine Nozzles}, KTH, 2002
--- Rao valid-region discussion.
\item R. Stark, \emph{Flow Separation in Rocket Nozzles --- An Overview}, 2009.
\item J. D. Anderson, \emph{Modern Compressible Flow} --- isentropic
relations, Prandtl--Meyer, MOC compatibility (\S11.4).
\item M. J. Zucrow \& J. D. Hoffman, \emph{Gas Dynamics}, Vol.\ 2, Ch.\ 17 ---
axisymmetric characteristics.
\item G. P. Sutton \& O. Biblarz, \emph{Rocket Propulsion Elements} ---
performance relations.
\end{itemize}

\begin{notebox}
\small Where any status statement in historical research notes conflicts with
the executable code or the current test suite, \textbf{the code and tests are
authoritative}. See \code{latex-report/raosim\_reference.pdf} for the extended
mathematical reference and \code{JAX\_DIFFERENTIABLE\_PLAN.md} for the
differentiable-solver development record.
\end{notebox}

\end{document}
